\section{Reanimated Husks}

Despite Reanimated Husks term can cover many enemies, they all share some common elements: they all have been alive once, but after their demise they were reanimated back to life as and undead puppets controlled by the Red Vines.

Due to the Red Vines affecting animated body, all of them receive Strength (and all skills, attributes and attacks affected by it) boost with the feature of \abilityVinesName{}:

\abilityVines{}

Since they are already dead, they also benefit from \abilityHuskName{} Feature:

\abilityHusk{}

As a result, the Reanimated Husk is a strong enemy focusing on brutal attacks, but it lacks finesse of a trained warrior - it cannot Graze and use its Focus abilities.

Since they are not fully alive, killing them doesn't need to end their existence. Some parts of the Husk's body can still contain enough Red Vines to continue fighting as \enemyReanimatedRemains{}. Those are severed limbs or other body parts that are puppeteered back into the fight.

In case of all living creatures, the bigger they are, the bigger the influence of the Red Vines. Both \abilityVinesName{} and \abilityHuskName{} scale with the size of the creature affected by them: \\

\begin{tabular}{c c c}
	Size 		& Scale & Ability Example \\ \hline
	Small 		& 1 & \abilityVinesName[1] \\
	Medium 		& 2 & \abilityVinesName[2] \\
	Large 		& 3 & \abilityVinesName[3] \\
	Huge 		& 4 & \abilityVinesName[4] \\
	Gargantuan 	& 5 & \abilityVinesName[5] \\
\end{tabular}

\tacticsBox{
	All Reanimated Husks attack whoever they can, and they never stop. They don't care about pain or death - they are already dead. All they want is to tear flesh, draw blood and spread the Red Vines' infestation.
}


